\documentclass{beamer}

\mode<presentation> {
\usetheme{Madrid}
}

\usepackage{graphicx}
\usepackage{booktabs}

% Beamer loads hyperref automatically. We configure link colors cleanly
\hypersetup{
    colorlinks=true,
    urlcolor=blue,
    linkcolor=blue
}

%----------------------------------------------------------------------------------------
%   TITLE PAGE
%----------------------------------------------------------------------------------------

\title[Maritime Cyber Range]{Simulating Cyberattacks on Autonomous Maritime Vessels}

\author{
    Vinaybaalaji PNS \quad (CB.SC.U4CYS24061) \\
    \medskip
    Department of TIFAC-CORE in Cyber Security \\
    Amrita Vishwa Vidyapeetham, Coimbatore \\
    \medskip
    3rd Year, B.Tech Cybersecurity
}

\institute[]{}
\date{\today}

\begin{document}

% Slide 1: Title Slide
\begin{frame}
\titlepage
\end{frame}

%----------------------------------------------------------------------------------------
%   SLIDES
%----------------------------------------------------------------------------------------

% Slide 2: Autonomous Vehicles Importance & Cyber Resilience
\section{Context}
\begin{frame}
\frametitle{\textbf{Autonomous Maritime Vessels and Cyber Resilience}}
\textbf{Role of Autonomous Vessels:}
\begin{itemize}
    \item Unmanned Surface Vessels (USVs) and Autonomous Underwater Vehicles (AUVs) are increasingly utilized for maritime patrol, ocean mapping, subsea inspection, and cargo transport.
    \item They handle long-duration and high-risk operations where deploying human crews is expensive or hazardous.
\end{itemize}

\vspace{0.4cm}
\textbf{Importance of Cyber Resilience:}
\begin{itemize}
    \item Unlike manned ships, autonomous vessels lack a crew on board to cross-check navigational anomalies or manually take control during a disruption.
    \item They rely entirely on wireless communication, satellite positioning, and acoustic telemetry.
    \item If these communication or sensor feeds are manipulated, the vessel blindly follows corrupted data, making cyber resilience critical to preventing collisions, mission aborts, or vessel loss.
\end{itemize}
\end{frame}

% Slide 3: Need for Testbeds, Regulations & Indian Landscape
\section{Regulatory Drivers}
\begin{frame}
\frametitle{\textbf{Pre-Sea Testing: Mandates \& National Context}}

\footnotesize
\textbf{Global Mandates Requiring Cyber Validation:}
\begin{itemize}
    \setlength\itemsep{0.1em}
    \item \textbf{IACS UR E26/E27 (2024):} Mandates cyber resilience testing for shipboard systems and onboard networks \href{https://iacs.org.uk}{\textbf{[1]}}.
    \item \textbf{IMO Res. MSC.428(98):} Requires cyber risk management in vessel safety management systems \href{https://www.imo.org}{\textbf{[2]}}.
    \item \textbf{ITU \& SOLAS Spectrum Laws:} Prohibits open-water radio jamming and GPS spoofing near civilian sea lanes.
\end{itemize}

\vspace{0.1cm}
\textbf{Indian Context \& Maritime Problem Statements:}
\begin{itemize}
    \setlength\itemsep{0.1em}
    \item DRDO NPOL's SPACE platform \href{https://www.pib.gov.in/PressReleasePage.aspx?PRID=2018151}{\textbf{[3]}} and NSTL basins evaluate acoustic physics, but lack active cyber testbeds.
    \item \textbf{MoD SIH PS 26247:} Real-time maritime threat detection and tactical intelligence fusion system \href{https://www.sih.gov.in/sih2026PS}{\textbf{[4]}}.
    \item \textbf{MoD SIH PS 26248:} Degraded communications and EW simulation for naval tactical command \href{https://www.sih.gov.in/sih2026PS}{\textbf{[5]}}.
    \item \textbf{Indian Navy NIIO / iDEX SPRINT:} Cyber-resilient autonomous naval craft (USVs/AUVs) and secure acoustic data links \href{https://idex.gov.in}{\textbf{[6]}}.
\end{itemize}

\vspace{0.1cm}
\tiny
\textbf{References:} 
\href{https://iacs.org.uk}{[1] IACS E26} \quad 
\href{https://www.imo.org}{[2] IMO 428} \quad 
\href{https://www.pib.gov.in/PressReleasePage.aspx?PRID=2018151}{[3] DRDO SPACE} \quad
\href{https://www.sih.gov.in/sih2026PS}{[4] SIH 26247} \quad
\href{https://www.sih.gov.in/sih2026PS}{[5] SIH 26248} \quad
\href{https://idex.gov.in}{[6] NIIO SPRINT}
\end{frame}

% Slide 4: Our Approach
\section{Our Approach}
\begin{frame}
\frametitle{\textbf{Our Approach: Simulation-Based Cyber Testing}}
\small
Our objective is to test cyberattacks across both surface (USV) and subsurface (AUV) craft in a controlled software simulation.

\vspace{0.2cm}
\textbf{Key Elements of the Approach:}
\begin{itemize}
    \setlength\itemsep{0.2em}
    \item \textbf{Cross-Domain Testing:} Evaluating attacks on both surface vessels (radio/GPS layer) and underwater vessels (acoustic communication layer).
    \item \textbf{Realistic Communications:} Modeling MAVLink 2 C2 command streams and telemetry between the ground control station and the vessel autopilot.
    \item \textbf{Indian Naval Test Sites:} Calibrated to realistic ocean conditions at Kochi, Mumbai, and Visakhapatnam.
    \item \textbf{Dual-Engine Ocean Dynamics:} Coupling Gazebo Sim (DART surface hydrodynamics) and Stonefish (6-DOF subsea added mass and damping).
    \item \textbf{Safe Experimentation:} Assessing navigation stability, EKF divergence, and boundary conditions without risking physical naval assets.
\end{itemize}
\end{frame}

% Slide 5: Prototype Demo (3D Visuals)
\section{Prototype Demo}
\begin{frame}
\frametitle{\textbf{Initial Prototype Demonstration: 3D Tactical Interface}}
\begin{columns}
    \column{0.5\textwidth}
    \centering
    \includegraphics[width=\linewidth]{godot_3d_interceptor_usv.png}
    \column{0.5\textwidth}
    \centering
    \includegraphics[width=\linewidth]{godot_3d_maya_auv.png}
\end{columns}

\vspace{0.3cm}
\footnotesize
\begin{itemize}
    \item \textbf{Left (Surface USV 3D Visual):} High-speed Fast Interceptor USV navigating in Godot 4 with dynamic ocean wave shaders, sun reflections, wake, and live HUD telemetry.
    \item \textbf{Right (Subsurface AUV 3D Visual):} Maya AUV operating submerged with procedural seabed shaders, marine snow particulate physics, and acoustic depth telemetry.
\end{itemize}
\end{frame}

% Slide 6: SITL Methodology (Orderly)
\section{Methodology}
\begin{frame}
\frametitle{\textbf{Simulation Methodology}}
Our simulation environment couples three core components in an orderly pipeline:

\vspace{0.3cm}
\begin{enumerate}
    \item \textbf{Dual-Physics Simulation (Gazebo \& Stonefish):} Simulates 6-DOF hydrodynamics, water resistance, buoyancy, added mass, and ocean current forces acting on the vessel hull.
    \vspace{0.2cm}
    \item \textbf{Autopilot Brain (ArduPilot / ArduSub SITL):} Runs production autopilot software in Software-in-the-Loop, executing real 24-state EKF3 sensor fusion and navigation algorithms.
    \vspace{0.2cm}
    \item \textbf{Communication Protocol (MAVLink 2):} Handles bidirectional telemetry and command routing between the autopilot and control station, matching real vessel network traffic.
\end{enumerate}
\end{frame}

% Slide 7: Vehicle Integration (Conditional)
\section{Extensibility}
\begin{frame}
\frametitle{\textbf{Extending to Other Vessels}}
\small
We can make this testbed vehicle-independent if the following inputs are available for a given platform:

\vspace{0.2cm}
\begin{enumerate}
    \setlength\itemsep{0.2em}
    \item \textbf{Autopilot Firmware / Binary:} The compiled control software or physical controller board used to execute guidance logic.
    \item \textbf{Sensor API:} An interface allowing the physics engine to stream simulated sensor data (depth, position, attitude) into the control logic.
    \item \textbf{Communication Protocol Specs:} Interface Control Documents (ICD) defining message formats so testing scripts can send recognizable commands.
    \item \textbf{Configuration Profile:} The parameter configuration of the vessel (speed limits, thruster mapping, and failsafe thresholds).
\end{enumerate}

\vspace{0.2cm}
\footnotesize
\textbf{Validated Digital Twins:} Verified across 4 autonomous platforms: Fast Interceptor USV, Coastal Survey Catamaran, Maya AUV, and AUV-150.
\end{frame}

% Slide 8: Attacks Tested & Real-World Fidelity
\section{Attack Vectors}
\begin{frame}
\frametitle{\textbf{Attacks Tested and Simulation Fidelity}}
\small
\textbf{The attacks we are testing:}
\begin{itemize}
    \setlength\itemsep{0.15em}
    \item \textbf{GPS Spoofing:} Injected coordinate shifts on surface vessels to alter navigation courses.
    \item \textbf{AIS Ghost Vessels \& Impersonation:} Forged NMEA AIVDM broadcast messages injecting phantom collision hazards.
    \item \textbf{Acoustic Position Spoofing:} Forged underwater positioning signals (USBL) where GPS cannot penetrate.
    \item \textbf{C2 Command Injection:} Forged MAVLink override packets attempting unauthorized mode changes and motor hijacking.
\end{itemize}

\vspace{0.2cm}
\textbf{Simulation Fidelity vs. Physical Reality:}
\begin{itemize}
    \setlength\itemsep{0.15em}
    \item \textbf{Protocol and Software Level:} Identical to reality. The autopilot firmware, packet headers, and state logic match production deployments.
    \item \textbf{Environmental Level:} Water density, currents, and regional sound speed are modeled; seabed acoustic reverberation and RF weather degradation remain simplified.
\end{itemize}
\end{frame}

% Slide 9: Benchmarks & Limitations (Updated with Live Tested Metrics)
\section{Benchmarks}
\begin{frame}
\frametitle{\textbf{Tested Simulation Benchmarks \& Operational Metrics}}

\centering
\resizebox{\linewidth}{!}{%
\begin{tabular}{lcc}
\toprule
\textbf{Parameter / Test Metric} & \textbf{Real Sea Environment} & \textbf{Tested Simulation Value} \\
\midrule
Seawater Density & $1020 - 1023\text{ kg/m}^3$ & $1020.5 - 1022.9\text{ kg/m}^3$ (Site-calibrated) \\
Sound Speed ($c$) & $1538 - 1543\text{ m/s}$ & $1538.7 - 1542.6\text{ m/s}$ (Mackenzie model) \\
Tidal Ocean Current & $0.30 - 0.70\text{ m/s}$ & $0.35 - 0.60\text{ m/s}$ (Dynamic drift vector) \\
Wave Height ($H_s$) \& Period & $1.0 - 1.2\text{ m}$ (Sea State 2--3) & $1.0 - 1.2\text{ m}$ ($T_p = 6.0 - 7.0\text{ s}$) \\
Subsea Acoustic Propagation & $\tau = d/c$, $4-6\text{ dB/km}$ & $\tau = 0.52\text{ s}$ ($800\text{ m}$), Thorp $25\text{ kHz}$ \\
Telemetry \& EKF Loop Rate & $10 - 50\text{ Hz}$ & $10.0\text{ Hz}$ MAVLink / $50\text{ Hz}$ EKF3 \\
\bottomrule
\end{tabular}%
}

\vspace{0.25cm}
\raggedright
\footnotesize
\textbf{Measured Attack Impact \& Detection Benchmarks (Live Audit):}
\begin{itemize}
    \setlength\itemsep{0.15em}
    \item \textbf{GNSS Spoofing:} Divergence of $10.0\text{ m} - 25.4\text{ m}$ with $11.2\text{ dB}$ $C/N_0$ front-end drop; detected by Page's CUSUM without prior signature knowledge.
    \item \textbf{AIS Ghost Insertion:} Forged phantom vessel triggered autonomous COLREGs Rule 14 collision avoidance with $30^\circ$ course deviation.
    \item \textbf{C2 Mode Takeover:} Forged MAVLink packets achieved 100\% unauthenticated authority seizure from GUIDED to MANUAL/HOLD.
    \item \textbf{Acoustic Position Spoofing:} Manipulated subsea USBL positioning within physical sound propagation limits ($\tau = 0.52\text{ s}$).
\end{itemize}
\end{frame}

\end{document}
